\section{Background and related work}
\label{sec:related}

\subsection{Evidence requests in claims handling}
\label{sec:context}

Insurers, legal-expenses insurers and public agencies handle claims in knowledge-intensive processes, in which the path a case takes depends on facts that emerge from documents and correspondence \cite{diciccio2015knowledge}. Deciding which documents to request is one recurring step. Its rules come from statutes, policy terms and internal guidelines, and its errors have direct costs: an unnecessary request delays the claim and burdens the claimant, a missed request lets the claim proceed without the evidence a payment or refusal needs, and a premature request asks for a document before the fact that makes it relevant is known. Claims handling is expected to be among the most automated parts of insurance \cite{eling2018impact,eling2022impact}, and process and information systems research asks how language models can act within explicit constraints, explain their actions and keep people in control \cite{dumas2023aiaugmented,vidgof2023large,kampik2025large,chau2025research,hevner2023research}. The request decision is a useful test case: frequent, rule-governed, consequential and checkable after the fact.

\subsection{Approaches}
\label{sec:approaches}

Table~\ref{tab:approaches} compares the approaches on five capabilities that our evaluation tests or uses: reading free-text claims, applying rules in code, holding unknown facts open, verifying judgments before use, and tracing each request to a rule.

\begin{table}[t]
\caption{What existing approaches provide for deciding evidence requests. \checkmark{} provides; $\circ$ partly; -- no. Columns: (1) reads free-text claims; (2) applies the rules in code; (3) holds unknown facts open instead of assuming them; (4) verifies judgments before use; (5) traces each request to a rule.}
\label{tab:approaches}
\small
\begin{tabularx}{\linewidth}{@{}Xccccc@{}}
\toprule
Approach & 1 & 2 & 3 & 4 & 5 \\
\midrule
Rule engines and rules as code \cite{omg2024dmn,merigoux2021catala,sergot1986british,mohun2020cracking} & -- & \checkmark & -- & -- & \checkmark \\
Case management and compliance checking \cite{omg2016cmmn,hull2011business,governatori2006compliance,hashmi2018compliance} & -- & \checkmark & $\circ$ & -- & $\circ$ \\
Policy or process in the prompt \cite{yao2025taubench,xiao2024flowbench,ye2025sopagent,li2026sopbench} & \checkmark & -- & -- & -- & -- \\
State- or graph-guided agents \cite{wu2024stateflow,lu2026proceduralgraphs,yu2026compile} & \checkmark & $\circ$ & -- & -- & $\circ$ \\
LLM rule application \cite{servantez2024chain,blairstanek2023statutory,holzenberger2020dataset} & \checkmark & -- & -- & -- & $\circ$ \\
Neuro-symbolic translation \cite{pan2023logic,olausson2023linc,ye2023satlm,gao2023pal} & \checkmark & \checkmark & -- & $\circ$ & $\circ$ \\
Legal rule programs with extracted facts \cite{holzenberger2023connecting,jurayj2026tax,kant2024equitable} & \checkmark & \checkmark & -- & -- & $\circ$ \\
Clarification and abstention \cite{li2024mediq,zhang2024askbeforeplan,kamath2020selective,wen2025abstention} & \checkmark & -- & \checkmark & -- & -- \\
Guardrails and verification \cite{chen2025shieldagent,xiang2025guardagent,dhuliawala2024chain} & \checkmark & $\circ$ & -- & \checkmark & $\circ$ \\
Evidence obligation planning \cite{kim2026evidenceobligation,taranukhin2026infogatherer} & \checkmark & $\circ$ & \checkmark & $\circ$ & \checkmark \\
\midrule
\casepath (this article) & \checkmark & \checkmark & \checkmark & \checkmark & \checkmark \\
\bottomrule
\end{tabularx}
\end{table}

\paragraph{Rules in code}
Executable rules have a long history in law and administration, from the British Nationality Act as a logic program \cite{sergot1986british} to Catala \cite{merigoux2021catala}, the Decision Model and Notation standard (DMN) \cite{omg2024dmn} and ``rules as code'' in government \cite{mohun2020cracking}. They apply rules deterministically and traceably, but take typed facts, not free text. A fact that is not supplied is an error or false: logic programs make this the closed-world assumption \cite{reiter1978closed}, and even DMN's three-valued expressions send a condition that is not true to the default branch instead of holding the decision open \cite{omg2024dmn}. Case management notations activate work when conditions are met \cite{omg2016cmmn,hull2011business}, and a pending sentry can wait for an event, which holds a fact open in part; compliance checking verifies processes against norms \cite{governatori2006compliance,hashmi2018compliance}. Both take conditions as structured events, and their logs trace an activation to the event that triggered it rather than to a quoted source.

\paragraph{Models that apply the rules}
The prevailing way to make an LLM agent follow a policy is to place it in the prompt; $\tau$-bench, FlowBench and procedure-following agents evaluate this design \cite{yao2025taubench,xiao2024flowbench,ye2025sopagent,li2026sopbench}. Other agents keep a state machine or procedural graph in code, which partly moves rule application into code, but the model still decides within each state \cite{wu2024stateflow,lu2026proceduralgraphs,yu2026compile}. Statutory-reasoning methods such as Chain of Logic decompose a rule into elements and let the model evaluate and recombine them \cite{servantez2024chain,blairstanek2023statutory,holzenberger2020dataset,guha2023legalbench}; their stated steps explain a decision in part. In all these designs the model judges the case and applies the rule in one generation, so neither decision is recorded separately, and repeated runs differ \cite{atil2025nondeterminism,blairstanek2025unstable}.

\paragraph{Models that formalise, code that executes}
Logic-LM, LINC, SatLM and program-aided reasoning translate a problem into a formal representation and let a solver compute the answer \cite{pan2023logic,olausson2023linc,ye2023satlm,gao2023pal}; some refine the formalisation when the solver reports an error, a partial check, and the formalisation shows in part how an answer was derived. The formalisation is produced anew for each problem, so the rules can change between cases, and the solver assumes every fact given or false. Legal systems translate statutes once and represent each case formally, by information extraction or by a language model \cite{holzenberger2023connecting,jurayj2026tax,kant2024equitable}; they compute an answer, such as a tax liability, from a complete case representation, with a program trace but no tie to quoted case text, rather than what to ask for next.

\paragraph{Asking, abstaining and verifying}
Clarification and selective-prediction methods let a model ask or abstain when it is uncertain \cite{rao2018learning,aliannejadi2019asking,kamath2020selective,wen2025abstention}, and interactive agents ask for missing information before acting \cite{li2024mediq,zhang2024askbeforeplan}. They decide from the model's confidence in its answer, not from whether an unresolved fact can change a rule's outcome. Guardrail agents check proposed actions against policies, some by compiling the policies into checkable rules that name the violated policy \cite{chen2025shieldagent,xiang2025guardagent}, and verification prompts let a model check its draft \cite{dhuliawala2024chain}; they filter an action after the fact. The closest precedent keeps a ledger of evidence obligations in code, planned by the model for each query, checked in part, held open until resolved and tied to their sources \cite{kim2026evidenceobligation}; information-gathering agents choose questions by expected evidential value \cite{taranukhin2026infogatherer}. Neither executes a shared process whose branches decide which obligations exist.

\subsection{The gap}
\label{sec:gap}

Each capability exists, but not together. Rule engines execute fixed rules traceably but cannot read a claim or leave a fact open; LLM agents read the claim but apply the rules themselves; neuro-symbolic pipelines execute a per-case formalisation and close the world; clarification methods ask, but not because a rule depends on the answer. Its ingredients are established: three-valued logic, supervaluation, rule execution, model-based extraction and verification. What is missing is their combination for this decision, an account of what each kind of error becomes under it, and evidence about whether it improves on letting the model apply the rules.
